Evidence-centric work takes the context window itself as the object of study: what the
retriever puts in front of the generator, in what order, and at what budget
(Appendix~\ref{app:evidence}). The stakes are higher here than in classical retrieval,
because the retrieved set is not a ranked list a user inspects but a prompt a model reads,
and a single misleading passage can move the answer, the stated confidence, and the citations
all at once.

\paragraph{Noise is not one condition.}
The most consequential distinction in this literature is between failing to fetch the right
evidence and failing to ignore the wrong evidence. \citet{irrelevantContext2023} separate the
two explicitly, and \citet{raat2024} show that the second is far harder than it looks: clearly
off-topic passages are discarded easily, whereas semantically close counterfactuals and hard
negatives are not. RGB makes the same point by construction, scoring noise robustness,
negative rejection and counterfactual handling as three separate abilities rather than one
\citep{rgb2023}. The practical consequence is that a single robustness-under-noise number is
uninterpretable on its own, because two papers reporting it may be stressing different
abilities.

\paragraph{Context construction has negative marginal utility.}
Adding passages is not monotone. \citet{cuconasu2024powerofnoise} find that topically related
but wrong passages hurt generation while purely random text can occasionally stabilise it, and
\citet{liu2024lostmiddle} show that the same passage helps or hurts depending on where it sits
in the prompt. Longer contexts do not resolve this: hard negatives and misleading content
continue to degrade reasoning as the window grows \citep{rrag039,rrag003}. Benchmarks such as
CUB \citep{rrag032} and UDCG \citep{rrag058} follow this observation to its conclusion and score
retrieval by downstream utility rather than by relevance, treating the retriever as an active
constructor of the model's input rather than as a passive supplier of documents.

\paragraph{Mitigations trade one failure for another.}
Evidence-side defences---filtering \citep{rrag044}, compression and reordering \citep{rrag039},
utility-aware reranking \citep{rrag058}, and noise-aware training \citep{raat2024}---all buy
precision with recall, so a defence that improves accuracy under distraction may lose answers
that a noisier pipeline would have found. Interventions on the query and index side, such as
rephrasing \citep{queryRobustness2025} or chunking and graph construction
\citep{rrag141,rrag051}, add failure surfaces of their own. None of these results transfers
across papers unless the corpus snapshot, retriever, and top-$k$ are stated, which is the
content of coordinate C1 in \S\ref{sec:coordinates}.
